\section{Discussion and Limitations}
\label{sec:discussion}

\subsection{Why a low evasion rate is a result, not a failed attack}

Some readers might view an Attack Success Rate (ASR) of 0.0--5.1\% as a failed
attack. In our setting, the opposite reading is more accurate. The attack
spends its full edit budget. It changes 13.3--13.6\% of characters and
6.2--9.2\% of words. It runs up to three attempts per sample. In total, it flips
40 decisions. What the attack does not do is guided search. It draws edits at
random instead of selecting the most damaging ones.

Section~\ref{sec:saliency} shows why search matters. We test synonym
replacements under the same 15\% word budget. When substitutions are ordered by
measured saliency (their observed impact on the model's decision), the flip rate
is 4.6\%. When substitutions are chosen at random, the flip rate is only 2.7\%.
In paired tests, saliency flips 33 messages that random edits miss, while
random edits flip only 7 that saliency misses. Thus, we should not claim these
models are robust. Instead, under random bounded edits they rarely flip; under a
guided search with the same budget, they flip about twice as often.

\subsection{Word edits cost fewer changes and flip more decisions}

Word-level changes flip more decisions than character-level noise (28 versus 12
flips). At the same time, word edits modify fewer units (6.2--9.2\% of words
versus 13.3--13.6\% of characters). The paired analysis in
\S\ref{sec:pairs} confirms this pattern. Word edits outperform character edits
in five of six cells, and the pooled difference is significant (26 versus 10,
$p=0.011$). The text similarity values explain this outcome. Word edits preserve
92--95\% of original words. Character edits leave only 30.3--32.8\% word overlap.
A homoglyph (a character that looks like another) alters the token
representation even when human readers see no change.

\subsection{Vulnerability is concentrated in specific words}

The saliency results clarify how random attacks behave. If an evasion succeeds
only when striking specific decisive words, a random draw from a small
dictionary rarely hits them. That explains the low 2.7\% baseline flip rate.

The saliency attack confirms this behavior. Most evasions require only two or
three targeted word substitutions. However, the saliency attack edits more words
on average (9.3\% versus 7.6\%, as shown in Table~\ref{tab:budget}). Normalized
by edit rate, its advantage is roughly 1.4 times rather than 1.7 times. Thus,
vulnerability is concentrated in specific tokens. This finding is useful for
defenders because targeted words can be filtered. It also warns against taking
random perturbation rates as proof of general robustness.

This effect varies across models. Phi-3.5-mini sits close to its decision
boundary. Its ASR roughly doubles under saliency search (11.7\% and 7.4\%
versus 5.1\% and 3.7\% under random word edits). Qwen2.5-3B is more accurate and
does not flip more often under saliency (2.4\% and 0.6\% versus 2.6\% and 3.2\%).
Its remaining errors cannot be reached by this candidate vocabulary.

\subsection{Two readouts of one model are not one measurement}

The saliency attack requires continuous scores, so it optimizes the predicted
probability $P(\text{malicious})$ over two answer tokens. The rest of this
paper evaluates the generated text decision. We re-checked every stored
saliency sample with text generation. The two channels agree on 941 of 1{,}001
rows. The disagreements are systematic. For Phi-3.5-mini, both channels agree on
all rows. For Qwen2.5-3B, the probability channel claims 45 flips that
generation does not confirm. For SmolLM2-1.7B, generation claims 15 flips that
the probability route misses.

This gap occurs because a two-token probability readout does not match free-form
decoding. When answer words split into sub-tokens, probability ratios can point
one way while greedy decoding points the other. We report final text decisions
throughout the paper. We view the divergence as a limitation of probability
probing in small language models.

\subsection{The transfer we observe is low}

We replayed every stored perturbation (character, word, and saliency) against
every other model. Very few perturbations transferred. On the applied subset,
only 51 of 5{,}514 caught messages (0.9\%) flipped off the diagonal.
Specifically, 9 flips came from 1{,}860 character edits, 18 from 1{,}847 word
edits, and 24 from 1{,}807 saliency edits. Guided edits transferred slightly
better than random edits. However, confidence intervals overlap, indicating a
broad trend rather than a sharp separation. In the classic flipped-source
matrix, only 5 of 31 cells are non-zero, and 27 of 31 cells hold fewer than ten
caught originals.

A 0.9\% rate indicates low observed transfer, not total immunity. Fifty-one
messages did cross model boundaries. Either these perturbations are
model-specific, or they are too weak to cross distant decision boundaries. In
either case, robustness evaluations on one model do not guarantee safety on
another.

\subsection{Defenses: character cleaning succeeds while re-thresholding fails}

Character normalization cleans visually confusable characters. It removes every
character-level flip (1.0\% to 0.0\% pooled). It leaves word and saliency
attacks largely unaffected (2.5\% to 2.0\%, and 7.5\% to 6.9\%). This confirms
our hypothesis: distinct perturbation types fail for distinct reasons.

Re-thresholding yields mixed results. We adjust the threshold $\tau$ to enforce a
5\% clean false-positive budget. This reduces flips on well-calibrated models.
However, on SmolLM2-1.7B (where clean scores cluster near 1.0), it increases
evasions: 0.0\% to 29.0\% on email and 0.0\% to 13.3\% on SMS. Lowering the
threshold places the decision cutoff inside perturbation noise. It also lowers
clean recall on SmolLM2-1.7B (100\% to 34.5\% and 23.0\%). For one task cell,
the budget cannot be met without rejecting all messages.

\subsection{Limitations}

\paragraph{Search strength.} Random perturbations are weak adversaries. The
saliency attack uses greedy first-order search over a fixed candidate set. It
is not a full implementation of PWWS or TextFooler, nor is it a white-box
gradient attack. Reported ASR values are lower bounds on attacker capabilities.

\paragraph{Transfer coverage.} The transfer matrix covers all three attacks.
Faithfulness is verified by self-replay: diagonal cells match the original
attack logs exactly. These tests demonstrate harness consistency on a single GPU
environment rather than universal cross-platform reproducibility.

\paragraph{Unmatched vocabulary in primary tables.} The saliency attack uses a
broader vocabulary than the random word attack. The ablation in
\S\ref{sec:ablation} isolates vocabulary, ordering, and oracle effects across
five arms. The reported contrasts in Table~\ref{tab:ablation} hold candidate
sets and query caps constant.

\paragraph{Defense operating point.} The defense probe evaluates a 5\%
false-positive budget using probability scores. Because thresholding requires
numeric scores, its baseline column is not directly comparable to text verdict
tables.

\paragraph{Saliency denominator.} The saliency attack excludes 43 originals that
its probability channel classified as benign (all on Qwen2.5-3B). Its denominator
is 1{,}001 messages, compared to 2{,}079 for random attacks. Nine additional
exclusions in the word attack correspond to identical text transformations (52
exclusions total).

\paragraph{Similarity metric.} We report word-level Jaccard similarity without
neural sentence embeddings. For character edits, Jaccard similarity is an
imperfect proxy, and actual edit rate is the cleaner measure.

\paragraph{Uncertainty scope.} Confidence intervals represent sample variation
over messages, not seed variability across runs.

\paragraph{Sample design.} Attacks target the correctly classified subset of an
800-row evaluation split. This provides an accurate conditional denominator but
does not represent out-of-distribution generalization.

\paragraph{Model coverage.} We evaluate three open models. Llama-3.2-3B-Instruct
was omitted due to access restrictions.

\paragraph{Task construct.} SMS smishing is evaluated on the UCI SMS Spam
dataset, where spam includes broader commercial promotions. Phishing email data
comes from public benchmark collections.

\paragraph{Missing experiments.} Paraphrase generation attacks and defense-aware
adaptive attacks were designed but not run.
